\documentclass[11pt]{article}

\usepackage{amsmath,amssymb,amsthm}
\usepackage[ruled,vlined,linesnumbered]{algorithm2e}
\usepackage{hyperref}
\usepackage{booktabs}
\usepackage{natbib}
\usepackage[margin=1in]{geometry}
\usepackage[expansion=false]{microtype}
\usepackage{graphicx}
\usepackage[section]{placeins}

\newtheorem{theorem}{Theorem}[section]
\newtheorem{lemma}[theorem]{Lemma}
\newtheorem{corollary}[theorem]{Corollary}
\newtheorem{proposition}[theorem]{Proposition}
\newtheorem{definition}[theorem]{Definition}
\newtheorem{remark}[theorem]{Remark}

\newcommand{\bigO}{\mathcal{O}}
\newcommand{\ip}[2]{\langle #1,\, #2 \rangle}

\title{Fast Tournament Equity Computation via \\
       Generating-Function Quadrature and FFT-Accelerated Subproduct Trees}

\author{Sam Rosenstrauch \\
        \texttt{sarose550@gmail.com}}

\date{July 21, 2026}

\begin{document}
\maketitle

\begin{abstract}
The Independent Chip Model (ICM) maps poker tournament chip stacks to
expected prize money, but exact computation requires summing over all
$n!$ elimination orderings. I present a deterministic algorithm that
reformulates ICM equity as a one-dimensional integral over
generating-function coefficients, evaluated by Gaussian quadrature
($Q = 256$ nodes, relative error ${<}\,1.6 \times 10^{-10}$). The
central challenge - computing leave-one-out polynomial products for all
$n$ players - is solved by an FFT-accelerated subproduct tree whose
propagation phase is the adjoint of its build phase, reducing cost from
$\bigO(nk)$ to $\bigO(n \log^2 k)$ per quadrature point. An empirical
lookup table (calibrated per device) dispatches
automatically between a batched linear engine and a hybrid block-tree
engine. I evaluate the CPU implementation in detail on Apple M3~Pro and
AMD Zen~4 (with AOCL-FFTW): the hybrid and linear engines compute exact
equities for up to 17,984 players in one second, single-threaded, on
Zen~4. A CUDA implementation for NVIDIA B200 extends this to roughly
1.5 million players in about a second.
\end{abstract}

% ======================================================================
\section{Introduction}\label{sec:intro}
% ======================================================================

Multi-table poker tournaments are structured competitions in which
players begin with equal chip stacks and are progressively eliminated as
they lose their chips. A payout structure $\pi = (\pi_0, \pi_1, \ldots, \pi_{k-1})$ specifies the prize awarded to each finishing position, with
$k \leq n$ positions receiving nonzero prizes and $\pi_0 \geq \pi_1 \geq
\cdots \geq \pi_{k-1} \geq 0$ (position $r+1$ receives payout $\pi_r$). The \emph{Independent Chip Model} (ICM)
provides a principled mapping from a player's chip stack to their
expected tournament equity - the expected prize money they will
receive~\citep{malmuth1987,harville1973}.

Despite its central role in tournament strategy, exact ICM computation
is fundamentally difficult because the number of possible elimination
orderings grows factorially in $n$. For $n = 20$ there are $2.4 \times
10^{18}$ orderings to consider; for $n = 30$, the count reaches
$2.7 \times 10^{32}$. Bitmask dynamic programming reduces this to
$\bigO(n \cdot 2^n)$ time by iterating over subsets of surviving
players~\citep{malmuth1987}, but the exponential state space limits
practical computation to roughly $n \leq 23$. Monte Carlo sampling
provides approximate solutions for larger fields~\citep{icmizer2024},
but error shrinks only as $\bigO(1/\sqrt{N})$, making high-precision
results prohibitively expensive.

The key observation that enables a better approach is that each player's
ICM equity can be expressed as a one-dimensional integral rather than an
exponential sum. This follows from the \emph{exponential clock framing},
where each player is assigned an independent exponential random variable
and players are eliminated in order of increasing firing time. Under
this framing, the combinatorial sum collapses into an integral that
depends only on the chip stacks.

The integral is evaluated by Gaussian quadrature with $Q = 256$ nodes,
achieving exponential convergence to double-precision accuracy - a choice
driven by worst-case stack ratios up to $10^9$, not just the well-behaved
uniform case (Section~\ref{sec:accuracy} justifies the specific value).
The remaining challenge is evaluating the integrand efficiently at each
quadrature point, which requires computing a degree-$k$ leave-one-out
polynomial product for each of $n$ players. I address this with an
FFT-accelerated binary subproduct tree whose propagation phase is the
adjoint of the build phase - a cross-correlation that distributes the
payout vector to all leaves simultaneously.

My contributions are:
\begin{enumerate}
  \item A rigorous derivation of ICM equities as coefficient extractions
    from a generating function, evaluated via numerical quadrature
    (Sections~\ref{sec:derivation}--\ref{sec:algorithm}).
  \item A hybrid engine combining a subproduct tree for inter-block
    computation with direct polynomial arithmetic for intra-block
    computation, with correctness proved via the adjoint relationship
    between convolution and cross-correlation (Section~\ref{sec:algorithm}).
  \item An empirical crossover table, calibrated offline via direct
    timing on the target hardware, verified to select the correct
    engine on every tested $(n,k)$ cell in serial mode on M3~Pro,
    replacing earlier analytical cost models (Section~\ref{sec:dispatch}).
  \item Platform-specific implementations for ARM64 (Apple M3~Pro),
    x86-64 (AMD Zen~4 with AOCL-FFTW), and CUDA (NVIDIA B200 GPU),
    demonstrating the approach at scale (Sections~\ref{sec:experiments},
    \ref{sec:gpu}, and~\ref{sec:gpu-results}).
\end{enumerate}

% ======================================================================
\section{Related Work}\label{sec:related}
% ======================================================================

\paragraph{ICM and the Malmuth--Harville model.}
The probability model underlying ICM was introduced by
\citet{harville1973} for horse-racing handicapping and adapted to poker
tournaments by \citet{malmuth1987}. Under this model, the probability
that any surviving player is eliminated next is proportional to their
chip stack. The standard exact algorithm uses bitmask dynamic programming
over all $2^n$ subsets of surviving players, running in
$\bigO(n \cdot 2^n)$ time. For large fields, Monte Carlo sampling is the
prevailing approach, with error converging as $\bigO(1/\sqrt{N})$.

\paragraph{Exponential clock formulation.}
The equivalence between the Malmuth--Harville elimination process and
independent exponential clocks was noted by Tysen Streib in a 2011
forum post~\citep{streib2011}, who used it to develop an efficient Monte
Carlo sampling method. The key insight is that assigning each player an
$\mathrm{Exp}(S_j)$ random variable and sorting by realized value
produces an elimination order whose distribution matches the
sequential-elimination model exactly, while requiring only $n$ random
draws and one sort per sampled tournament.

\paragraph{Polynomial arithmetic and subproduct trees.}
The subproduct tree (also called a \emph{subresultant tree} or
\emph{product tree}) is a standard tool in computer algebra for
multi-point evaluation and interpolation.
\citet{vonzurgathen2013} provide a comprehensive treatment, including the
$\bigO(n \log^2 n)$ algorithms for computing the product of $n$ linear
factors and evaluating a polynomial at $n$ points. My use of the tree is
closest to the ``accumulating remainder'' formulation, but with the
critical difference that propagation computes cross-correlations (the
adjoint of convolution) rather than polynomial remainders. This adjoint
perspective, proved in Theorem~\ref{thm:tree-correct}, appears to be new
in this context.

\paragraph{FFT libraries and automatic tuning.}
All FFT-accelerated engines in this work build on the Cooley--Tukey
algorithm~\citep{cooley1965}, which reduces the DFT from $\bigO(d^2)$ to
$\bigO(d \log d)$. FFTW~\citep{frigo2005} introduced the idea of empirical
auto-tuning for FFT plans, generating machine-specific ``codelets'' at
install time.
AMD's AOCL-FFTW extends this with znver4-tuned AVX-512-specific codelets
(versus the generic upstream distribution)~\citep{aocl2024}. I use
AOCL-FFTW as the sole FFT backend on Zen~4 after a direct A/B measurement
against plain system FFTW3 (Section~\ref{sec:experiments}) showed a clean,
reproducible win at the kernel level with no sizes where the reverse held
; so no dual-library dispatch is warranted on this platform, unlike the
GPU's cuFFT/cuFFTDx split (Section~\ref{sec:gpu-results}). The broader
philosophy of empirically-tuned numerical libraries traces to
FFTW's \texttt{PATIENT} planner~\citep{frigo2005} and
SPIRAL~\citep{pueschel2005}.  The approach of searching a parameter
space empirically per architecture, as done here for the
$(n,k)$-indexed crossover and block-size tables, is closest in
spirit to the register-blocking search in the BeBOP/Sparsity
system~\citep{demmel2004}, which times candidate register-blocking
and unrolling parameters for sparse matrix kernels.

\paragraph{Roofline performance models.}
The roofline model of \citet{williams2009} bounds achievable performance
by the minimum of peak compute and memory bandwidth, parameterized by
arithmetic intensity.  An earlier version of the engine dispatch
(Section~\ref{sec:dispatch}) used a roofline model with harmonic-mean
bandwidth blending across cache levels; this was superseded by the
current empirical crossover table after the analytical model proved
unable to match measured crossover points on real hardware.

\paragraph{GPU FFT and CUDA Graphs.}
NVIDIA's cuFFT library supports batched FFT execution, and the cuFFTDx
extensions~\citep{nvidia2024cufftdx} enable fused device-side FFTs
that eliminate intermediate memory round-trips. CUDA
Graphs~\citep{nvidia2024cudagraphs} amortize kernel launch overhead by
capturing and replaying entire execution graphs, which is essential for
the dozens of kernel launches in a single tree traversal.

% ======================================================================
\section{Mathematical Derivation}\label{sec:derivation}
% ======================================================================

This section derives the integral representation that is the foundation
of the algorithm. All of the key ideas - the exponential clock, the
generating function, and the one-dimensional integral - are presented
here, with supporting proofs where they provide insight.

\subsection{The Malmuth--Harville Model and Exponential Clocks}

\begin{definition}[Tournament instance]\label{def:tournament}
A \emph{tournament instance} is a tuple $(n, S, \pi)$ where
$n \in \mathbb{N}$ is the number of players, $S = (S_1, \ldots, S_n)
\in \mathbb{R}_{>0}^n$ is the vector of chip stacks with $S_i > 0$ for all
$i$, and $\pi = (\pi_0, \ldots, \pi_{k-1}) \in \mathbb{R}_{\geq 0}^k$
is the payout vector with $k \leq n$ and $\pi_0 \geq \pi_1 \geq \cdots
\geq \pi_{k-1} \geq 0$ (0-indexed: position $r+1$ receives payout $\pi_r$).
\end{definition}

\begin{definition}[Malmuth--Harville model]\label{def:mh}
Under the Malmuth--Harville (MH) model, the probability that player $i$ is
the first to be eliminated from a set $T \subseteq \{1, \ldots, n\}$ of
surviving players is:
\begin{equation}\label{eq:mh-elim}
  \Pr[i \text{ elim.\ first} \mid T] = \frac{S_i}{\sum_{j \in T} S_j}\,.
\end{equation}
Elimination proceeds sequentially: after the first player is eliminated, the
model is applied recursively to the remaining set $T \setminus \{i\}$ with
unchanged stacks.
\end{definition}

The MH model has an equivalent continuous-time formulation that is central
to my approach. Assign each player~$j$ an independent exponential random
variable $T_j \sim \mathrm{Exp}(S_j)$ with density $S_j e^{-S_j t}$ for
$t \geq 0$. Players are eliminated in order of increasing $T_j$ - smallest
$T_j$ is eliminated first, etc.

Why does this reproduce the MH model? For any two players $i, j$, the
probability that $i$ has a smaller clock than $j$ is:
\[
  \Pr[T_i < T_j] = \frac{S_i}{S_i + S_j}.
\]
More generally, for any set of surviving players, the memoryless property
of the exponential distribution guarantees that the conditional probability
of any surviving player being eliminated next is proportional to their
stack, recovering Equation~\ref{eq:mh-elim} exactly.

This means we can sample an entire elimination order in one shot: draw $n$
independent exponentials, sort by their realized values, and the resulting
order is distributed according to the MH model. This is the basis for
efficient Monte Carlo ICM methods, but it also enables something stronger.

\subsection{The Generating Function}

Fix player $i$ and condition on $T_i = t$. For each other player $j \neq i$,
we have, independently:
\begin{itemize}
  \item $j$ finishes \emph{before} $i$ iff $T_j < t$, which happens with
    probability $b_j(t) = \Pr[T_j < t] = 1 - e^{-S_j t}$.
  \item $j$ finishes \emph{after} $i$ with probability
    $a_j(t) = \Pr[T_j > t] = e^{-S_j t}$.
\end{itemize}

Because the clocks are independent, the probability of any specific pattern
of ``who finishes before $i$'' is the product of the individual
probabilities. This product structure suggests a generating function.
Define:
\begin{equation}\label{eq:gen-func-def}
  Q_i(x; t) = \prod_{j \neq i} \bigl(a_j(t) + b_j(t) \cdot x\bigr).
\end{equation}

Expanding this product, each factor contributes either $a_j(t)$ (player $j$
finishes after $i$, contributing $x^0$) or $b_j(t) \cdot x$ (player $j$
finishes before $i$, contributing $x^1$). The coefficient of $x^r$ in
$Q_i(x; t)$ is therefore exactly the probability that exactly $r$ of the
other players finish before $i$, conditioned on $T_i = t$:
\begin{equation}\label{eq:coeff-meaning}
  [x^r]\, Q_i(x; t) = \Pr[i \text{ finishes in position } r+1 \mid T_i = t].
\end{equation}

\subsection{The Integral Representation}

To uncondition, we integrate against the density of $T_i$:
\begin{equation}\label{eq:position-prob}
  \Pr[i \text{ finishes in position } r+1]
  = \int_0^\infty S_i e^{-S_i t} \cdot [x^r]\, Q_i(x; t) \, dt.
\end{equation}

Now change variables: let $v = e^{-t}$, so $dv = -e^{-t}\,dt = -v\,dt$ and
$dt = -dv/v$. Substituting $a_j(t) = e^{-S_j t} = v^{S_j}$ and
$b_j(t) = 1 - v^{S_j}$:
\begin{equation}\label{eq:gf-v}
  Q_i(x; v) = \prod_{j \neq i} \bigl(v^{S_j} + (1 - v^{S_j}) x\bigr).
\end{equation}

As $t$ goes from $0$ to $\infty$, $v$ goes from $1$ to $0$. The minus sign
from $dt = -dv/v$ cancels the bound flip, yielding:
\begin{equation}\label{eq:position-prob-v}
  \Pr[i \text{ finishes in position } r+1]
  = \int_0^1 S_i \, v^{S_i - 1} \cdot [x^r]\, Q_i(x; v) \, dv.
\end{equation}

Multiplying by the payout $\pi_r$ and summing over positions:
\begin{equation}\label{eq:equity-integral}
  E_i = \int_0^1 S_i \, v^{S_i - 1} \,
    \underbrace{\sum_{r=0}^{k-1} \pi_r \cdot [x^r]\, Q_i(x; v)}_{%
      \textstyle =: \Psi_i(v)} \, dv.
\end{equation}

The inner sum $\Psi_i(v) = \sum_{r=0}^{k-1} \pi_r \cdot [x^r] Q_i(x; v)$
is the dot product of the payout vector with the coefficient vector of
$Q_i$ - exactly the quantity the subproduct tree will compute efficiently.

\subsection{Numerical Quadrature}

The integral in Equation~\ref{eq:equity-integral} is over $v \in (0, 1)$.
For numerical integration, I apply the change of variables
$v = \Phi(y)$ where $\Phi$ is the standard normal CDF. This maps
$(0, 1)$ to $\mathbb{R}$ and produces an integrand that decays like
$e^{-y^2/2}$ as $|y| \to \infty$, making it ideal for the trapezoidal
rule~\citep{trefethen2014}:
\begin{equation}\label{eq:quadrature}
  E_i \approx \sum_{q=1}^{Q} w_q \cdot S_i \cdot v_q^{S_i - 1}
    \cdot \Psi_i(v_q),
\end{equation}
where $(v_q, w_q)_{q=1}^Q$ are the quadrature nodes and weights.

This rule (which I call \emph{erfc-trap}) achieves exponential convergence:
the error shrinks roughly as $e^{-cQ}$ for some $c > 0$ depending on the
strip of analyticity of the integrand. With $Q = 256$, relative error is
below $5 \times 10^{-12}$ on uniform stacks and below $1.6 \times
10^{-10}$ even at the $10^9$ stack-ratio bound (see
Section~\ref{sec:accuracy} for a detailed comparison against
tanh-sinh quadrature).

% ======================================================================
\section{Algorithm}\label{sec:algorithm}
% ======================================================================

The quadrature rule reduces ICM computation to evaluating
$\Psi_i(v_q) = \sum_{r=0}^{k-1} \pi_r \cdot [x^r] Q_i(x; v_q)$ for
all $i = 1, \ldots, n$ at each of $Q$ quadrature points. Computing
$Q_i(x; v_q) = \prod_{j \neq i} (a_j + b_j x)$ naively - one player at a
time - would cost $\bigO(n^2 k)$ per quadrature point. This section
presents three engines that compute all $n$ leave-one-out products
simultaneously, with cost-model-driven automatic dispatch.

I will work throughout with truncated polynomials in the vector space
$\mathbb{R}[x]_{<k} \cong \mathbb{R}^k$ of polynomials of degree $<k$,
equipped with the Euclidean inner product
$\ip{f}{g} = \sum_{m=0}^{k-1} f_m g_m$, so that
$\Psi_i(v) = \ip{\pi}{Q_i}$.

\subsection{Engine 1: Batched Linear Forward-Backward}\label{sec:linear}

The linear engine computes all $n$ inner products in $\bigO(nk)$ time using
a forward-backward decomposition. It is optimal when $k$ is small.

\begin{algorithm}[H]
\caption{Linear Forward-Backward}\label{alg:linear}
\DontPrintSemicolon
\SetAlgoLined
\KwIn{$n$ players, values $a_1, \ldots, a_n$, 0-indexed payout $\pi \in \mathbb{R}^k$}
\KwOut{$\Psi_1, \ldots, \Psi_n$}
\BlankLine
\tcp{Forward: prefix products}
$g^{(0)} \leftarrow \pi$\;
\For{$j = 1, \ldots, n$}{
  \For{$m = 0, \ldots, k-2$}{
    $g^{(j)}_m \leftarrow a_j \cdot g^{(j-1)}_m + (1-a_j) \cdot g^{(j-1)}_{m+1}$\;
  }
  $g^{(j)}_{k-1} \leftarrow a_j \cdot g^{(j-1)}_{k-1}$\;
}
\BlankLine
\tcp{Backward: suffix products, extract inner products}
$(R_0, \ldots, R_{k-1}) \leftarrow (1, 0, \ldots, 0)$\;
\For{$j = n, n-1, \ldots, 1$}{
  $\Psi_j \leftarrow \sum_{m=0}^{k-1} g^{(j-1)}_m \cdot R_m$\;
  \For{$m = k-1, k-2, \ldots, 1$}{
    $R_m \leftarrow a_j \cdot R_m + (1-a_j) \cdot R_{m-1}$\;
  }
  $R_0 \leftarrow a_j \cdot R_0$\;
}
\Return{$\Psi_1, \ldots, \Psi_n$}
\end{algorithm}

The forward pass computes prefix products $g^{(j)} = \pi * P_1 * \cdots
* P_j \bmod x^k$, where $P_j(x) = a_j + (1-a_j)x$. The backward pass
constructs suffix products $R = P_{j+1} * \cdots * P_n \bmod x^k$
incrementally, extracting $\Psi_j = \ip{g^{(j-1)}}{R}$ at each step.
The inner product and suffix update are fused into a single loop because
the backward iteration over $m$ reads $R_m$ for the dot product before
overwriting it.

\paragraph{Batched quadrature and SIMD.}
To exploit SIMD parallelism, I interleave $B_Q = 8$ quadrature points in the
data layout: $a_{\text{batch}}[j \cdot B_Q + q_i]$ stores the $q_i$-th
quadrature value for player~$j$, ensuring sequential access in the inner
loops. $B_Q = 8$ wins on all platforms: on Zen~4 it maps to one native
512-bit AVX-512 vector per player, and on Apple M3~Pro it maps to two
NEON FMA operations spread across 4 FMA ports.

\paragraph{L2-aware checkpointing.}
When $n \cdot k \cdot B_Q \cdot 8$ bytes exceeds L2 cache, the forward
pass saves checkpoints every $C = \lfloor L_2 / (k \cdot B_Q \cdot 8)
\rfloor$ players. The backward pass recomputes each segment from its
checkpoint before extracting equities, trading $\bigO(n/C)$ recomputation
passes for $\bigO(C \cdot k)$ memory. This is critical not only for the
serial engine but also for the parallel (OpenMP) path, where omitting it
forces DRAM bandwidth instead of L2 and causes a severe regression.

\subsection{Engine 2: FFT-Accelerated Subproduct Tree}\label{sec:tree}

For large $k$, the linear engine's $\bigO(nk)$ cost becomes expensive.
The tree engine replaces it with a $\bigO(n \log^2 k)$ approach based on
a binary subproduct tree.

\begin{definition}[Binary subproduct tree]
Given $n$ linear factors $P_j(x) = a_j + (1-a_j)x$ for $j = 1, \ldots, n$,
the binary subproduct tree has $L = \lceil \log_2 n \rceil + 1$ levels.
Level~0 contains the $n$ leaves $P_j$. Level~$\ell$ contains
$\lceil n / 2^\ell \rceil$ nodes, each being the product of its two
children (truncated to degree $k$):
\[
  P^{(\ell)}_i = P^{(\ell-1)}_{2i} \cdot P^{(\ell-1)}_{2i+1} \bmod x^k.
\]
\end{definition}

\begin{algorithm}[H]
\caption{Tree Build and Propagate}\label{alg:tree}
\DontPrintSemicolon
\SetAlgoLined
\KwIn{Leaf polynomials $P^{(0)}_1, \ldots, P^{(0)}_n$,
  0-indexed payout $\pi$, degree $k$}
\KwOut{$\Psi_1, \ldots, \Psi_n$}
\BlankLine
\tcp{Build (bottom-up)}
\For{$\ell = 1, \ldots, L-2$}{
  \For{each node $i$ at level $\ell$ \textbf{in parallel}}{
    $P^{(\ell)}_i \leftarrow P^{(\ell-1)}_{2i} \cdot P^{(\ell-1)}_{2i+1} \bmod x^k$
  }
}
\BlankLine
\tcp{Propagate (top-down)}
$g^{(L-1)} \leftarrow \pi$\;
\For{$\ell = L-1, L-2, \ldots, 1$}{
  \For{each node $i$ at level $\ell$ \textbf{in parallel}}{
    $g^{(\ell-1)}_{2i} \leftarrow g^{(\ell)}_i \star P^{(\ell-1)}_{2i+1}$
    \tcp*{Left child via right sibling}
    $g^{(\ell-1)}_{2i+1} \leftarrow g^{(\ell)}_i \star P^{(\ell-1)}_{2i}$
    \tcp*{Right child via left sibling}
  }
}
\Return{$g^{(0)}_1[0], \ldots, g^{(0)}_n[0]$}
\end{algorithm}

\subsection{The Adjoint: Why Propagation Works}

The correctness of the tree engine follows from a single algebraic fact:
the adjoint of truncated convolution is cross-correlation.

\begin{definition}
For $h \in \mathbb{R}[x]_{<k}$, let $T_h : \mathbb{R}[x]_{<k} \to
\mathbb{R}[x]_{<k}$ denote truncated convolution by $h$:
$T_h(f) = f * h$, i.e.\ $(T_h(f))_m = \sum_{i+j=m} f_i h_j$. Define
cross-correlation by $f$ and $h$ as
$(f \star h)_j = \sum_{m=0}^{k-1-j} h_m f_{m+j}$.
\end{definition}

\begin{lemma}[Correlation is the adjoint of convolution]\label{lem:adj-conv}
For all $f, g, h \in \mathbb{R}[x]_{<k}$,
$\ip{T_h(f)}{g} = \ip{f}{g \star h}$; equivalently, $T_h^{\ast} =
{\cdot} \star h$.
\end{lemma}
\begin{proof}
$T_h$ has a lower-triangular Toeplitz matrix representation
$M_{ij} = h_{i-j}$ (with $h_\ell = 0$ for $\ell < 0$ or $\ell \ge k$).
Then $\ip{Mf}{g} = f^\top M^\top g = \ip{f}{M^\top g}$, and
$(M^\top g)_i = \sum_{j=0}^{k-1} h_{j-i} g_j =
\sum_{m=0}^{k-1-i} h_m g_{m+i} = (g \star h)_i$.
\end{proof}

The build phase constructs the full product $R = P_1 * \cdots * P_n$ via a
binary tree of truncated convolutions. Reading from root to leaf $j$, the
product decomposes as $R = Q_j * P_j \bmod x^k$, where $Q_j$ is the product
of all \emph{sibling} polynomials on the path from leaf $j$ to the root.

The propagation phase computes $\ip{\pi}{Q_j}$ for all $j$ by
applying the adjoint of the build phase's linear map, with $\pi$ as
the seed. At each internal node, the build phase computes
$(L, R) \mapsto L * R \bmod x^k$. Holding the sibling fixed and viewing
this as the linear map $T_S$ (where $S$ denotes the sibling polynomial)
in the descendant of interest,
the adjoint $T_S^{\ast}$ is $g \mapsto g \star S$
- exactly what Algorithm~\ref{alg:tree} implements at each level.

\begin{theorem}[Correctness of the tree engine]\label{thm:tree-correct}
After Algorithm~\ref{alg:tree} completes, $g^{(0)}_j[0] = \Psi_j$ for all
$j = 1, \ldots, n$.
\end{theorem}
\begin{proof}
We establish the invariant by downward induction on $\ell$. Let
$\mathrm{leaves}(i,\ell)$ be the leaf indices in the subtree rooted at node
$i$ at level $\ell$, and let $C_j^{(\ell)}$ be the set of all leaves
\emph{outside} this subtree. The invariant is:
\begin{equation}\label{eq:invariant}
  g^{(\ell)}_i = \pi \star
    \Bigl(\prod_{j \in C_j^{(\ell)}} P_j \bmod x^k \Bigr).
\end{equation}

\textbf{Base} ($\ell = L{-}1$, root): $C_j^{(L-1)} = \emptyset$, and
$\pi \star 1 = \pi = g^{(L-1)}$.

\textbf{Inductive step}: Assume the invariant holds at level $\ell$ for
node $i$. The left child receives
$g^{(\ell-1)}_{2i} = g^{(\ell)}_i \star P^{(\ell-1)}_{2i+1}$. For any
test vector $f$:
\begin{align*}
  \ip{f}{g^{(\ell-1)}_{2i}}
  &= \ip{f}{(\pi \star R_C) \star P_{2i+1}} \\
  &= \ip{f * P_{2i+1}}{\pi \star R_C}
    \quad\text{(Lemma~\ref{lem:adj-conv})} \\
  &= \ip{(f * P_{2i+1}) * R_C}{\pi}
    \quad\text{(Lemma~\ref{lem:adj-conv})} \\
  &= \ip{f * (P_{2i+1} * R_C)}{\pi}
    \quad\text{(associativity)} \\
  &= \ip{f}{\pi \star (P_{2i+1} * R_C)}
    \quad\text{(Lemma~\ref{lem:adj-conv})}.
\end{align*}
Since this holds for all $f$, the invariant holds for the left child with
$C_j^{(\ell-1)} = C_j^{(\ell)} \cup \mathrm{leaves}(2i+1, \ell-1)$. The
right child is symmetric. At level 0, $C_j^{(0)} = [n] \setminus \{j\}$,
so $g^{(0)}_j = \pi \star Q_j$ and
$g^{(0)}_j[0] = \ip{\pi}{Q_j} = \Psi_j$.
\end{proof}

\begin{remark}
The propagation phase is an instance of reverse-mode automatic
differentiation applied to the build phase~\citep{griewank2008}. Writing
$T_\ell$ for the convolution map $T_h$ of Lemma~\ref{lem:adj-conv}
instantiated with $h$ equal to level $\ell$'s sibling polynomial, the
build computes $T_L \circ \cdots \circ T_1$, and propagation computes
$T_1^{\ast} \circ \cdots \circ T_L^{\ast}$ applied to $\pi$. A
standard result in reverse-mode AD states that the adjoint of a linear
program has the same arithmetic complexity, giving complexity parity
between build and propagation for free.
\end{remark}

\paragraph{Per-level FFT and paired cached correlate.}
At each tree level, I compare calibrated FFT costs against schoolbook
multiplication to decide whether to use FFT. FFT sizes are chosen from
all 7-smooth numbers ($2^a 3^b 5^c 7^d$) to minimize total cost including
wrap correction. During propagation, I share the forward FFT of the parent
$g$-vector across both child correlates and cache child frequency-domain
representations from the build phase, halving the FFT count at each level.

This sharing is possible because correlation and convolution differ only
by a conjugation in the frequency domain: for real-valued $f, h$, the
discrete Fourier transform satisfies $\widehat{f \star h} = \hat f \cdot
\overline{\hat h}$, whereas ordinary convolution satisfies $\widehat{f * h}
= \hat f \cdot \hat h$ (no conjugate). Conjugating a transformed vector is
$\bigO(d)$ - an elementwise sign flip on the imaginary part - versus
$\bigO(d\log d)$ for a fresh transform. Two consequences
follow directly. First, each child polynomial's forward FFT, already
computed once during the build phase to convolve it with its sibling
(Algorithm~\ref{alg:tree}, build step), needs only to be conjugated - not
retransformed - to serve as that child's correlation kernel during
propagation; this is the ``cached'' frequency-domain representation.
Second, propagation at a node correlates the \emph{same} parent $g$-vector
against both children's (cached, conjugated) spectra
($g^{(\ell)}_i \star P^{(\ell-1)}_{2i+1}$ and
$g^{(\ell)}_i \star P^{(\ell-1)}_{2i}$), so a single forward FFT of
$g^{(\ell)}_i$ - the ``paired'' share - suffices for both child updates
instead of transforming $g^{(\ell)}_i$ separately for each, halving the
FFT count at this level relative to a naive implementation.

\subsection{Engine 3: Hybrid Block-Divide}\label{sec:hybrid}

The hybrid engine partitions the $n$ players into $\lceil n/B \rceil$
blocks of size $B$ and combines direct $\bigO(B^2)$ within-block
computation with the tree engine for inter-block structure.

\begin{algorithm}[H]
\caption{Hybrid Block-Divide}\label{alg:hybrid}
\DontPrintSemicolon
\SetAlgoLined
\KwIn{$n$ players, values $a_1, \ldots, a_n$, 0-indexed payout $\pi$, degree $k$, block size $B$}
\KwOut{$\Psi_1, \ldots, \Psi_n$}
\BlankLine
\tcp{Step 1: Build block products}
\For{each block $b = 0, 1, \ldots, \lceil n/B \rceil - 1$}{
  $P_b(x) \leftarrow \prod_{j \in \text{block } b} (a_j + (1-a_j)x)$ \tcp*{Sequential, $\bigO(B^2)$}
}
\BlankLine
\tcp{Step 2: Inter-block tree (Algorithm~\ref{alg:tree}) with block leaves}
Build subproduct tree from block products $P_0, \ldots, P_{\lceil n/B \rceil - 1}$\;
Propagate $\pi$ to obtain block-level $g$-vectors $g_b$\;
\BlankLine
\tcp{Step 3: Within-block divide}
\For{each block $b$}{
  \For{each player $j$ in block $b$}{
    Compute $Q_j^{(b)} \leftarrow P_b(x) / (a_j + (1-a_j)x)$
      via stable bidirectional synthetic division\;
    $\Psi_j \leftarrow \sum_{m=0}^{\min(k,B)-1} g_{b,m} \cdot Q_{j,m}^{(b)}$\;
  }
}
\Return{$\Psi_1, \ldots, \Psi_n$}
\end{algorithm}

\paragraph{Bidirectional division stability.}
Write the complete block product as $P_b(x) = \sum_{m=0}^{B} P_m x^m$ and
the quotient to be recovered as $Q_j^{(b)}(x) = \sum_{m=0}^{B-1} Q_m x^m$,
so that $P_b = (a_j + (1-a_j)x) \cdot Q_j^{(b)}$. Matching coefficients on
both sides gives two exact recurrences for $Q$, depending on which
direction the coefficients are read in:

\emph{Forward} (increasing $m$): $Q_0 = P_0 / a_j$, and for
$m = 1, \ldots, B-1$,
\[
  Q_m = \frac{P_m - (1-a_j)\,Q_{m-1}}{a_j}
      = \frac{P_m}{a_j} - \frac{1-a_j}{a_j}\, Q_{m-1}.
\]

\emph{Backward} (decreasing $m$): $Q_{B-1} = P_B / (1-a_j)$, and for
$m = B-2, \ldots, 0$,
\[
  Q_m = \frac{P_{m+1} - a_j\,Q_{m+1}}{1-a_j}
      = \frac{P_{m+1}}{1-a_j} - \frac{a_j}{1-a_j}\, Q_{m+1}.
\]

Both compute the same exact quotient in exact arithmetic; they differ only
in the coefficient $c$ multiplying the running term each step - forward's
$c_{\text{fwd}} = -(1-a_j)/a_j$, backward's $c_{\text{bwd}} = -a_j/(1-a_j)$ -
which governs how a rounding error at step $m$ is amplified by step $m+1$.

\begin{lemma}[Stability of bidirectional division]\label{lem:division-stability}
Choosing the forward recurrence when $a_j > 1/2$ and the backward
recurrence when $a_j \le 1/2$ gives per-step error amplification
$|c| \le 1$ for all $a_j \in (0,1)$, with $|c| < 1$ strictly except at the
boundary $a_j = 1/2$.
\end{lemma}
\begin{proof}
$|c_{\text{fwd}}| = (1-a_j)/a_j < 1 \iff a_j > 1/2$, and
$|c_{\text{bwd}}| = a_j/(1-a_j) \le 1 \iff a_j \le 1/2$; the two conditions
are exhaustive and disjoint except at $a_j = 1/2$ itself, where both give
$|c| = 1$.
\end{proof}

Lemma~\ref{lem:division-stability} proves $|c| \le 1$ rigorously: when the
stable branch is chosen, the per-step amplification factor never exceeds
unity. However, a tight analytical bound on the \emph{accumulated} rounding
error as a function of $B$, connecting $|c|^B$ near the worst-case
$a_j \approx 0.5$ to FP64's machine epsilon ($\approx 2.2 \times 10^{-16}$,
roughly 15--16 significant decimal digits), was never derived in closed
form. The division at each step introduces a factor of $1/a_j$ (or
$1/(1-a_j)$), which approaches~$2$ when $|c| \approx 1$, so the full
error growth depends on more than $|c|^B$ alone.

Empirically, correctness was verified up to $B = 256$ on AMD Zen~4:
bidirectional division with $B \in \{8, 16, 24, 32, 48, 64, 96, 128, 192,
256\}$ was tested against the $V_1$ exact reference across all
\texttt{bench\_grid verify} sizes and stack distributions, including the
adversarial near-$a_j = 0.5$ regime the lemma flags as riskiest. All
tests passed with no error growth across the entire range (standard:
$\max\ 1.18 \times 10^{-12}$ vs.\ $5 \times 10^{-12}$ tolerance;
adversarial: $\max\ 1.54 \times 10^{-10}$ vs.\ $10^{-9}$ tolerance).
Separately, an exhaustive block-size sweep on the same hardware found
that $B = 32$ is typically the fastest choice, and larger $B$ values are
never faster. The production candidate list
$\{8, 16, 24, 32, 48, 64\}$ therefore caps at $64$ for performance
reasons, not because of a hard correctness wall.

\paragraph{Rejected CPU optimizations.} Several plausible optimizations
were tried on the CPU engines and rejected after benchmarking. Karatsuba
multiplication for the block build was consistently slower: on FMA-capable
hardware, schoolbook's \texttt{c[i+j] += a[i]*b[j]} is already one FMA per
term, so Karatsuba's trade of multiplies for additions saves nothing - both
cost one cycle. FFTW's NEON codelets for FP64 were slower than Apple's
vDSP on Apple Silicon, which is why vDSP is dispatched for FFT execution
instead of FFTW at supported sizes on that platform
(Section~\ref{sec:experiments}).

\subsection{Engine Dispatch: Empirical Lookup Table}\label{sec:dispatch}

Engine dispatch must decide, for each $(n,k)$ pair, whether to use the
linear engine (memory-bound, $\bigO(nk)$) or the hybrid engine
(compute-bound in the FFT phases, $\bigO(n \log^2 k)$). An earlier
FMA-counting roofline model with harmonic-mean bandwidth blending correctly
dispatched roughly 30/44 test cases on Zen~4: it had no knowledge of the
cache hierarchy and could not distinguish L2-resident from DRAM-resident
working sets. A later roofline model~\citep{williams2009} improved this to
42/44 (95.5\%) by modeling cache-level bandwidths explicitly, but the two
remaining misses occurred where the engines were within 3--9\% of each
other, and individual constants in the formula were validated against real
embedded execution only to find that the aggregate comparison still did not
match the true measured crossover on either M3~Pro or Zen~4 hardware.

The replacement is an empirical crossover table: $N = 6$ calibrated $n$
values with corresponding crossover $k$ values where linear and hybrid
engines have equal measured runtime, determined offline by binary search on
real timing (median of 7 repetitions, $Q = 256$). At runtime,
\texttt{select\_engine}$(n,k)$ performs log-linear interpolation between
the two bracketing $n$ values in the \texttt{crossover\_n[]} table to
obtain the crossover $k$ for the requested $n$, then dispatches to linear
if $k$ is below the crossover and hybrid otherwise.  This is the same
structural precedent as LAPACK's \texttt{ILAENV} \texttt{ISPEC=3} (the
\texttt{NX} parameter)~\citep{anderson1999}: a problem-size crossover
between two algorithms, measured empirically per machine, consulted as a
cheap runtime threshold with no live racing in production.

The block size $B$ for the hybrid engine is selected by the same principle:
\texttt{select\_best\_B}$(n,k)$ performs a 2D nearest-neighbor lookup over
a calibrated $(n,k,B)$ grid (2466 points for M3~Pro and 1944 for Zen~4), choosing the
empirically fastest $B \in \{8, 16, 24, 32, 48, 64\}$.  Unlike the
crossover $k$, $B$ is a discrete choice with no meaningful interpolation
between adjacent values.  Typical selections are $B = 32$ on both M3~Pro
and Zen~4, with occasional $B = 24$ or $B = 48$ at specific $(n,k)$
combinations.

Both tables (\texttt{crossover\_n[]}/\texttt{crossover\_k[]} and
\texttt{bselect\_n[]}/\texttt{bselect\_k[]}/\texttt{bselect\_B[]}) are
generated offline by \texttt{tools/calibrate\_crossover.c} and
\texttt{tools/calibrate\_best\_b.c} respectively and stored in the
per-device \texttt{fft\_config.h}.  This approach follows the broader
autotuning tradition of measuring rather than modeling: FFTW's
\texttt{PATIENT} planner times candidate FFT plans directly instead of
relying on an indirect cost heuristic~\citep{frigo2005}, and, for the
specific mechanism of searching a parameter space empirically per
architecture, is closest to the register-blocking search in the
BeBOP/Sparsity system~\citep{demmel2004}, though that work searches
register-blocking/unrolling parameters for one fixed small kernel size
rather than building a size-indexed lookup table as done here.

\paragraph{Validation.}
\texttt{tools/validate\_dispatch.c} directly compares
\texttt{select\_engine}$(n,k)$'s choice against the true faster engine
(measured head to head, adaptive median-of-$N$ convergence) over the same
$(n,k)$ grid as Table~\ref{tab:serial} -- 7 values of $n$ from $1{,}024$
to $65{,}536$ crossed with $k \in \{10, 50, 100, n/4, n/2, n\}$, 42 cells.
On M3~Pro, serial mode, dispatch was correct on all 40 cells completed in
this evaluation run (the two largest, $n=65{,}536$ at $k=n/2$ and $k=n$,
were not measured due to a runtime constraint in this environment, not a
dispatch failure); the surrounding cells at that same $n$, including
$k=n/4$, were all correct.

% ======================================================================
\section{Theoretical Analysis}\label{sec:theory}
% ======================================================================

\begin{theorem}[Complexity of the tree engine]
The tree engine (Algorithm~\ref{alg:tree}) computes all $n$ inner products
in $\bigO(n \log^2 k)$ operations per quadrature point.
\end{theorem}
\begin{proof}
Let $T_b(n)$ be the build cost for $n$ leaves with truncation degree $k$.
At the root, two subtrees of $n/2$ leaves are combined by truncated
polynomial multiplication costing $\bigO(\min(n, k) \log \min(n, k))$ via
FFT, giving $T_b(n) = 2T_b(n/2) + \bigO(\min(n, k) \log \min(n, k))$.

\textbf{Case 1} ($n \le k$): $T_b(n) = 2T_b(n/2) + \bigO(n \log n)$,
solved by the Master theorem to $T_b(n) = \bigO(n \log^2 n)$.

\textbf{Case 2} ($n > k$): all subproducts are capped at degree $k$, so
merge cost is $\bigO(k \log k)$ regardless of subtree size. The recursion
bottom is reached when subtree size $n/2^\ell \le k$, i.e.\ at level
$\ell = \log_2(n/k)$. Below this level, the tree has at most $n/k$
independent subtrees of size $\le k$, each costing $\bigO(k \log^2 k)$
by Case~1; a total of $(n/k) \cdot \bigO(k \log^2 k) = \bigO(n \log^2 k)$.
Above saturation, there are $\log_2(n/k)$ levels each with
$\le n/(k \cdot 2^{\ell})$ merges at $\bigO(k \log k)$ each, summing to
$\sum_{\ell=1}^{\log_2(n/k)} \frac{n}{k \cdot 2^\ell} \cdot \bigO(k \log k)
= \bigO(n \log k \cdot \sum_{\ell} 2^{-\ell}) = \bigO(n \log k)$.
Adding the two contributions gives $\bigO(n \log^2 k + n \log k) =
\bigO(n \log^2 k)$.

In both cases, $T_b(n) = \bigO(n \log^2 k)$. Propagation has the same
complexity by adjoint parity (Remark after Theorem~\ref{thm:tree-correct}).
\end{proof}

\begin{theorem}[Complexity of the hybrid engine]\label{thm:hybrid-complexity}
The hybrid engine with block size $B$ computes all $n$ inner products in
$\bigO(nB + n\log^2 k)$ operations per quadrature point.
\end{theorem}
\begin{proof}
Block build: $n/B$ blocks at $\bigO(B^2)$ each = $\bigO(nB)$.
Within-block divide: $\bigO(B)$ per player = $\bigO(nB)$.

For the inter-block tree, let $m = n/B$ leaves, each a degree-$B$
polynomial. The tree has $\log_2 m$ merge levels above the leaves.
Define the \emph{saturation level} $s = \log_2(k/B)$: this is the number
of levels where subtree degree remains below $k$, i.e.\ $B \cdot 2^\ell
< k$ for $\ell \le s$ and capped at~$k$ thereafter.

\noindent\textit{Below saturation} ($\ell = 1,\ldots,s$). There are
$m/2^\ell$ nodes at level~$\ell$, each merging two degree-$(B \cdot
2^{\ell-1})$ polynomials into one of degree $B \cdot 2^\ell$, costing
$B \cdot 2^\ell \log_2(B \cdot 2^\ell)$ per merge. Per-level cost:
$(m/2^\ell) \cdot B \cdot 2^\ell \log_2(B \cdot 2^\ell)
= n \log_2(B \cdot 2^\ell)$. Summing over $\ell$:
\[
  C_{\text{below}} = n\sum_{\ell=1}^{s} \log_2(B \cdot 2^\ell)
  = n\sum_{\ell=1}^{s} (\log_2 B + \ell)
  = n\Bigl(s\log_2 B + \frac{s(s+1)}{2}\Bigr).
\]
Since $s = \log_2 k - \log_2 B$, the sum telescopes:
\begin{align*}
  s\log_2 B + \frac{s(s+1)}{2}
  &= \frac{s}{2}\bigl(2\log_2 B + s + 1\bigr) \\
  &= \frac{s}{2}\bigl(\log_2 B + (\log_2 B + s) + 1\bigr) \\
  &= \frac{s}{2}\bigl(\log_2 B + \log_2 k + 1\bigr)
   = \frac{s}{2}\log_2(2kB).
\end{align*}
Thus $C_{\text{below}} = \frac{n}{2}\,s\,\log_2(2kB)$.

\noindent\textit{At and above saturation} ($\ell = s{+}1,\ldots,\log_2 m$).
All merges cost $k\log_2 k$ each. The number of nodes across these levels
is a geometric series:
\[
  \sum_{\ell=s+1}^{\log_2 m} \frac{m}{2^\ell}
  = m\bigl(2^{-s} - 2^{-\log_2 m}\bigr)
  = \frac{n}{B}\Bigl(\frac{B}{k} - \frac{B}{n}\Bigr)
  = \frac{n}{k} - 1.
\]
Hence $C_{\text{above}} = (n/k - 1) \cdot k\log_2 k = (n - k)\log_2 k$.

Total tree cost:
\begin{equation}\label{eq:tree-cost-exact}
  C_{\text{tree}}(n,B,k)
  = \frac{n}{2}\,\log_2\!\Bigl(\frac{k}{B}\Bigr)\,
    \log_2(2kB) \;+\; (n - k)\log_2 k.
\end{equation}

At $B = k$ we have $s = 0$ and the first term vanishes, leaving
$C_{\text{tree}} = (n-k)\log_2 k$, which is exactly the ``all-saturated''
cost of a tree whose leaves already meet the degree cap. At $B = 1$
we have $s = \log_2 k$, giving
$C_{\text{tree}} = \frac{n}{2}\log_2 k \cdot \log_2(2k) + (n-k)\log_2 k
= \frac{n}{2}\log_2^2 k + O(n\log k)$.
For any $B \le k$, Equation~\ref{eq:tree-cost-exact} is
$\Theta(n\log^2 k)$: the first term is $\Theta(n\log^2(k/B))$ and the
second is $\Theta(n\log k)$; together they are $\Theta(n\log^2 k)$
because $\log(k/B) \le \log k$.
\end{proof}

\begin{corollary}[Blocking does not improve the asymptotic order]\label{cor:blocking-constant}
The total hybrid cost $T(B) = nB + C_{\text{tree}}(n,B,k)$ is minimized
at $B = \Theta(1)$, and increasing $B$ beyond a small constant strictly
increases the total.
\end{corollary}
\begin{proof}
With Equation~\ref{eq:tree-cost-exact} and $s = \log_2(k/B)$,
the change when doubling $B$ is
$\Delta T = T(2B) - T(B) = n(2B - B) + \frac{n}{2}\bigl[(s-1)\log_2(4kB)
- s\log_2(2kB)\bigr]$.  Substituting $\log_2(4kB) = \log_2(2kB) + 1$ and
simplifying:
\begin{align*}
  \Delta T &= nB + \frac{n}{2}\bigl[(s-1)(\log_2(2kB) + 1)
              - s\log_2(2kB)\bigr] \\
           &= nB + \frac{n}{2}\bigl[s - \log_2(2kB) - 1\bigr] \\
           &= nB + \frac{n}{2}\bigl[(\log_2 k - \log_2 B)
              - (\log_2 k + \log_2 B + 1) - 1\bigr] \\
           &= nB - n(\log_2 B + 1)
            = n\bigl(B - \log_2 B - 1\bigr).
\end{align*}
For $B = 1$: $\Delta T = n(1 - 0 - 1) = 0$, so $T(2) = T(1)$.
For $B = 2$: $\Delta T = n(2 - 1 - 1) = 0$, so $T(4) = T(2)$.
For $B \ge 4$: $f(B) = B - \log_2 B - 1$ has derivative
$f'(B) = 1 - 1/(B\ln 2) > 0$ for $B > 1/\ln 2 \approx 1.44$, and
$f(4) = 4 - 2 - 1 = 1 > 0$.  Hence $\Delta T > 0$ for all $B \ge 4$,
and $T(B)$ is strictly increasing on $B \in \{4,8,16,\ldots\}$.

Thus $T(1) = T(2) = T(4) < T(8) < T(16) < \cdots$, and the minimum is
attained at $B \in \{1,2,4\}$.  In particular, $B = \Theta(1)$.
\end{proof}

\paragraph{Theory versus practice.}
Corollary~\ref{cor:blocking-constant} shows that blocking provides no
asymptotic speedup: the tree phase is $\Theta(n\log^2 k)$ regardless of
$B$, and the $nB$ block-build term grows faster than the tree-phase
savings for any $B \ge 4$.  The optimal $B$ is therefore determined
entirely by hardware constant factors: cache-line width, SIMD lane
count, division latency; not by $n$ or $k$.  This is exactly what the
empirical block-size sweep finds (Section~\ref{sec:experiments}): the
wall-clock-optimal $B$ on Zen~4 is $B = 32$, constant across the entire
tested $(n,k)$ grid, with no detectable dependence on $k$ over the range
$k \in [10^2, 10^5]$.  Had blocking provided a $\Theta(\log k)$
asymptotic benefit, the optimal $B$ would grow with $k$; the fact that it
does not is an independent experimental confirmation of
Corollary~\ref{cor:blocking-constant}.

Numerical stability of the hybrid engine's within-block division is
established separately, alongside the division recurrence itself, in
Section~\ref{sec:hybrid} (Lemma~\ref{lem:division-stability}).

% ======================================================================
\section{CPU Experimental Results}\label{sec:experiments}
% ======================================================================

% All figures referenced below live in the sibling ICM repository:
% \graphicspath{{../ICM/results/}}

\subsection{Experimental Setup}

I evaluate the CPU implementation on two platforms (GPU hardware and setup
are described separately in Section~\ref{sec:gpu-results}). All experiments
use $Q = 256$ quadrature points, uniform chip stacks, and report the median
of 5 repetitions. Correctness is verified against closed-form $V_1$ and
$V_2$ equities (see Appendix) and cross-checked between engines.

\begin{itemize}
  \item \textbf{Apple M3~Pro:} 6 performance + 6 efficiency cores (12
    logical), ARM64 with NEON (2 FP64/vector), 4 FMA ports per core.
    FFT: FFTW3 PATIENT wisdom calibrated on this specific M3~Pro unit,
    with Apple vDSP dispatch at 33 supported sizes.
  \item \textbf{AMD Ryzen 9 7950X (Zen~4):} 16 homogeneous cores,
    AVX-512 (8 FP64/vector), 2 FMA units per core, 1\,MB per-core L2,
    32\,MB shared L3. FFT: AOCL-FFTW, znver4-tuned, sole backend (no
    dual-library dispatch; see Related Work). Memory on this machine runs
    at 3600\,MT/s rather than its DIMMs' 5600\,MT/s rating, an AM5
    two-DIMMs-per-channel electrical limit rather than a configuration
    fault. Its effect on these results is not a flat scaling factor: the
    compute-bound linear engine is essentially unaffected, while the
    FFT-heavy hybrid engine is memory-bandwidth-bound and correspondingly
    slower at large $k$. The Zen~4 hybrid columns below are therefore a
    bandwidth-limited lower bound for that microarchitecture, not its
    ceiling.
\end{itemize}

\subsection{CPU Performance}

Tables~\ref{tab:serial} and~\ref{tab:parallel} report execution times
across a grid of $(n, k)$ pairs. All numbers are from real measurements on
the target hardware (\texttt{bench\_grid}; M3~Pro July 24, 2026, Zen~4
July 27, 2026); none are estimated or borrowed. The best engine per cell is shown, with
dispatch decisions made automatically by the empirical crossover table.

\begin{table}[t]
\centering
\caption{Single-threaded execution time (ms), $Q = 256$, uniform stacks.
  Best engine per cell: L = linear, H = hybrid, T = tree.}
\label{tab:serial}
\resizebox{\textwidth}{!}{%
\begin{tabular}{@{}l rrrrrr rrrrrr@{}}
\toprule
& \multicolumn{6}{c}{\textbf{Apple M3 Pro}} &
  \multicolumn{6}{c}{\textbf{AMD Zen 4 (AOCL-FFTW)}} \\
\cmidrule(lr){2-7} \cmidrule(lr){8-13}
$n$ & $k{=}10$ & $k{=}50$ & $k{=}100$ & $k{=}n/4$ & $k{=}n/2$ & $k{=}n$
    & $k{=}10$ & $k{=}50$ & $k{=}100$ & $k{=}n/4$ & $k{=}n/2$ & $k{=}n$ \\
\midrule
1024  & 1.72(L) & 7.07(L) & 13.0(L) & 17.9(H) & 21.0(H) & 28.6(H)
      & 1.46(L) & 3.48(L) & 7.32(L) & 19.1(H) & 22.3(H) & 24.1(H) \\
2048  & 4.10(L) & 14.1(L) & 26.1(L) & 44.4(H) & 51.8(H) & 56.3(H)
      & 3.67(L) & 8.21(L) & 15.0(L) & 48.4(H) & 53.7(H) & 57.2(H) \\
4096  & 8.13(L) & 28.3(L) & 52.0(L) & 108(H)  & 123(H)  & 141(H)
      & 7.39(L) & 17.8(L) & 25.8(L) & 115(H)  & 128(H)  & 139(H) \\
8192  & 16.3(L) & 59.7(L) & 104(L)  & 291(H)  & 374(T)  & 407(T)
      & 13.4(L) & 34.3(L) & 53.6(L) & 277(H)  & 315(H)  & 341(H) \\
16384 & 32.4(L) & 113(L)  & 208(L)  & 784(T)  & 788(H)  & 967(T)
      & 29.6(L) & 56.9(L) & 117(L)  & 667(H)  & 743(H)  & 810(H) \\
32768 & 64.9(L) & 226(L)  & 416(L)  & 1670(H) & 1920(H) & 2100(H)
      & 62.3(L) & 131(L)  & 218(L)  & 1570(H) & 1790(H) & 1950(H) \\
65536 & 130(L)  & 452(L)  & 831(L)  & 4030(H) & 4670(H) & 5610(H)
      & 118(L)  & 266(L)  & 411(L)  & 4110(H) & 4750(H) & 5110(H) \\
\bottomrule
\end{tabular}%
}
\end{table}

\begin{table}[t]
\centering
\caption{Parallel execution time (ms), $Q = 256$, uniform stacks.
  M3 Pro: 12 threads (6P+6E). Zen 4: 16 threads.
  Best engine per cell: L = linear, H = hybrid, T = tree.}
\label{tab:parallel}
\resizebox{\textwidth}{!}{%
\begin{tabular}{@{}l rrrrrr rrrrrr@{}}
\toprule
& \multicolumn{6}{c}{\textbf{Apple M3 Pro (12 threads)}} &
  \multicolumn{6}{c}{\textbf{AMD Zen 4 (16 threads)}} \\
\cmidrule(lr){2-7} \cmidrule(lr){8-13}
$n$ & $k{=}10$ & $k{=}50$ & $k{=}100$ & $k{=}n/4$ & $k{=}n/2$ & $k{=}n$
    & $k{=}10$ & $k{=}50$ & $k{=}100$ & $k{=}n/4$ & $k{=}n/2$ & $k{=}n$ \\
\midrule
1024  & 0.325(L) & 1.28(L)  & 2.07(H)  & 2.42(H)  & 2.82(H)  & 3.81(H)
      & 0.121(L) & 0.359(L) & 1.21(L)  & 1.41(H)  & 1.66(H)  & 1.78(H) \\
4096  & 1.40(L)  & 4.90(L)  & 8.82(L)  & 14.8(H)  & 16.6(H)  & 18.9(H)
      & 0.651(L) & 1.48(L)  & 2.35(L)  & 8.40(H)  & 9.36(H)  & 10.1(H) \\
8192  & 2.71(L)  & 9.63(L)  & 16.8(H)  & 39.8(H)  & 50.6(T)  & 55.0(T)
      & 1.23(L)  & 2.84(L)  & 4.63(L)  & 20.4(H)  & 24.2(H)  & 28.1(H) \\
16384 & 5.40(L)  & 18.9(L)  & 33.9(H)  & 108(T)   & 108(H)   & 135(T)
      & 2.37(L)  & 6.43(L)  & 9.69(L)  & 88.1(H)  & 112(H)   & 142(H) \\
32768 & 10.7(L)  & 37.8(L)  & 68.1(L)  & 231(H)   & 265(H)   & 295(H)
      & 6.04(L)  & 13.6(L)  & 23.6(L)  & 280(H)   & 314(H)   & 363(H) \\
65536 & 21.1(L)  & 78.9(L)  & 140(H)   & 584(H)   & 682(H)   & 811(H)
      & 21.0(L)  & 30.4(L)  & 48.4(L)  & 627(H)   & 729(H)   & 901(H) \\
\bottomrule
\end{tabular}%
}
\end{table}

\paragraph{Platform comparison.}
Zen~4 benefits from AVX-512's 8-wide FP64 FMA (versus NEON's 2-wide) in the
linear engine, and from AOCL-FFTW's znver4-tuned AVX-512 codelets in the
hybrid engine.

A direct A/B comparison of AOCL-FFTW versus the upstream (plain) FFTW
distribution was conducted on the same Zen~4 machine, at the same FFT size
($d = 32{,}768$), with the same \texttt{FFTW\_PATIENT} flag, under the
\texttt{performance} cpufreq governor. AOCL-FFTW completed the per-pipeline
FFT work with a 20--25\% advantage over plain FFTW, reproducible across
repeated runs, with no calibrated size where the reverse held. AOCL-FFTW is therefore used as the sole FFT
backend on Zen~4; no dual-library dispatch is needed on this platform
(contrast the GPU's cuFFT/cuFFTDx tier selection, Section~\ref{sec:gpu-results},
where dispatch genuinely depends on problem size).

M3~Pro benefits from Apple vDSP's interleaved DFT dispatch at 33 supported
FFT sizes (10--18\% faster than FFTW at those sizes) and from the
integration of NEON SIMD throughout the linear engine's batched path
($B_Q = 8$ quadrature points interleaved).

\paragraph{Parallel scaling asymmetry.}
The two engines scale differently on 16 Zen~4 cores. The reference box was
redeployed 2026-07-27 in a 2-DIMMs-per-channel (\emph{2DPC}) configuration,
an AMD AM5 electrical limit that caps RAM at 3600\,MT/s against a higher
DIMM rating; the earlier \emph{1DPC} configuration ran at full rated
bandwidth. This is kept as two datasets rather than one, because the
effect is not a flat percentage: linear/schoolbook timings are nearly
unchanged between the two ($\sim$0.97--1.0$\times$) while hybrid/FFT
timings are 40--65\% slower under 2DPC at large $k$. Figure~\ref{fig:1dpc-2dpc}
shows the two eras' one-second contours side by side. Restricted to
$(n,k)$ cells where the same engine wins in both the serial and parallel
grids (isolating the scaling effect from a shifting engine choice), the
linear engine's median parallel speedup is $9.7\times$ on the 2DPC
reference box and $9.8\times$ on 1DPC, while the hybrid engine's is
$13.1\times$ and $12.6\times$ respectively -- hybrid scales somewhat
better than linear across the range measured. This is not uniform across
$n$: at $n = 65{,}536$ specifically (2DPC), linear's speedup ranges
$5.6$--$8.75\times$ and hybrid's $5.67$--$6.56\times$, comparable to each
other rather than showing the gap the medians suggest.
The per-core L2 capacity argument for the linear engine's checkpointing
still holds (Zen~4: $16 \times 131.5 \approx 2{,}104$\,GB/s aggregate
streaming bandwidth if every thread stayed L2-resident, Table~\ref{tab:constants}),
but does not by itself explain the magnitude gap between the two
engines; identifying the mechanism is open follow-up work.

This asymmetry creates a sharp cliff at the linear-to-hybrid transition,
which on Zen~4 falls between $k = 100$ and $k = 200$: the parallel linear
engine at $k = 100$ handles 906{,}259 players within one second, but the
parallel hybrid engine at $k = 200$ handles only 195{,}356 - a
$4.6\times$ drop that reflects the bandwidth regime change, not an
algorithmic deficiency. The serial contour shows the same transition at the
same $k$ and in the same direction (117{,}264 at $k = 100$ falling to
45{,}085 at $k = 200$), but a shallower $2.6\times$ drop, since a single
thread does not saturate the DRAM bus.

\begin{figure}[t]
  \centering
  \includegraphics[width=0.6\textwidth]{../ICM/results/zen4_1dpc_vs_2dpc.png}
  \caption{Zen~4 one-second contour, 1DPC (full memory bandwidth) vs.\ 2DPC
    (the standing reference configuration, 3600\,MT/s ceiling), serial,
    $Q = 256$. The two curves diverge at large $k$, where the hybrid
    engine's FFT phase is bandwidth-bound; they nearly coincide at small
    $k$, where the linear engine dominates and is largely
    bandwidth-insensitive.}
  \label{fig:1dpc-2dpc}
\end{figure}

\begin{figure}[t]
  \centering
  \includegraphics[width=0.48\textwidth]{../ICM/results/contour_1s_m3pro.png}
  \includegraphics[width=0.48\textwidth]{../ICM/results/contour_1s.png}
  \caption{One-second contour: maximum $n$ achievable within 1\,s as a
    function of $k$. Left: Apple M3~Pro. Right: AMD Zen~4. Solid line:
    serial; dashed line: parallel. The cliff between $k = 100$ and
    $k = 200$ (Zen~4) marks the linear-to-hybrid engine transition.}
  \label{fig:contour}
\end{figure}

\begin{figure}[t]
  \centering
  \includegraphics[width=0.48\textwidth]{../ICM/results/parallel_speedup_m3pro.png}
  \includegraphics[width=0.48\textwidth]{../ICM/results/parallel_speedup.png}
  \caption{Parallel speedup at the one-second boundary. Left: M3~Pro
    (12 threads, 6P+6E). Right: Zen~4 (16 threads).}
  \label{fig:speedup}
\end{figure}

\begin{figure}[t]
  \centering
  \includegraphics[width=0.48\textwidth]{../ICM/results/engine_dispatch_m3pro.png}
  \includegraphics[width=0.48\textwidth]{../ICM/results/engine_dispatch.png}
  \caption{Engine dispatch decisions at the one-second boundary.
    Purple: linear engine. Green: hybrid engine. Left: M3~Pro.
    Right: Zen~4.}
  \label{fig:dispatch}
\end{figure}

\begin{figure}[t]
  \centering
  \includegraphics[width=0.48\textwidth]{../ICM/results/runtime_vs_n_cpu_m3pro.png}
  \includegraphics[width=0.48\textwidth]{../ICM/results/runtime_vs_n_cpu.png}
  \caption{Runtime vs.\ $n$ at fixed $k$, single-threaded. Left: Apple
    M3~Pro. Right: AMD Zen~4. The CPU analogue of the GPU runtime plot in
    Figure~\ref{fig:gpu-contour}.}
  \label{fig:runtime-cpu}
\end{figure}

\subsection{Quadrature Accuracy}\label{sec:accuracy}

The erfc-trap quadrature rule - trapezoidal rule after the change of
variables $v = \Phi(y)$ - achieves exponential convergence for smooth
integrands~\citep{trefethen2014}. I compare it against tanh-sinh
(double-exponential) quadrature~\citep{mori2001} across four stack
distributions:
\begin{itemize}
  \item \textbf{Uniform:} $S_i = 100$ for all $i$ (ratio 1).
  \item \textbf{Adversarial:} $S_1 = 10{,}000$, $S_i = 1$ for $i > 1$
    (ratio $10^4$).
  \item \textbf{Geometric:} $S_i = 2^{i-1}$, giving $S_{\max}/S_{\min}
    \approx 5 \times 10^5$ at $n = 20$.
  \item \textbf{Extreme adversarial:} $S_1 = 10^9$, $S_i = 1$ for $i > 1$
    (ratio $10^9$). This is the practical worst case.
\end{itemize}

For each configuration I measure the maximum relative error across all $n$
players against closed-form $V_2$ equities. Table~\ref{tab:accuracy}
reports results for $n = 20$, the hardest case.

\begin{table}[t]
\centering
\caption{Maximum relative error vs.\ quadrature order $Q$ for $n = 20$.
  Gauss = erfc-trap; TS = tanh-sinh. Machine epsilon (FP64) is
  $2.2 \times 10^{-16}$.}
\label{tab:accuracy}
\resizebox{\textwidth}{!}{%
\begin{tabular}{@{}r c c c c c c c c@{}}
\toprule
& \multicolumn{2}{c}{Uniform} & \multicolumn{2}{c}{Adversarial}
  & \multicolumn{2}{c}{Geometric} & \multicolumn{2}{c}{Extreme ($10^9$)}\\
\cmidrule(lr){2-3} \cmidrule(lr){4-5} \cmidrule(lr){6-7} \cmidrule(lr){8-9}
$Q$ & Gauss & TS & Gauss & TS & Gauss & TS & Gauss & TS \\
\midrule
4    & $3.9$                & $2.5{\times}10^{-1}$  & $1.0{\times}10^{0}$  & $1.2{\times}10^{1}$  & $5.0$                & $1.1{\times}10^{1}$  & $1.0{\times}10^{0}$  & $1.0{\times}10^{0}$ \\
8    & $6.2{\times}10^{-3}$ & $1.0$                 & $2.3{\times}10^{0}$  & $9.8{\times}10^{-1}$ & $2.4$                & $1.0{\times}10^{0}$  & $9.9{\times}10^{-1}$ & $1.0{\times}10^{0}$ \\
16   & $1.6{\times}10^{-1}$ & $6.6{\times}10^{-1}$  & $1.4{\times}10^{-2}$ & $1.6{\times}10^{0}$  & $7.7{\times}10^{-1}$ & $1.4{\times}10^{0}$  & $1.2{\times}10^{0}$  & $4.1{\times}10^{0}$ \\
32   & $4.3{\times}10^{-3}$ & $8.5{\times}10^{-4}$  & $5.2{\times}10^{-2}$ & $3.0{\times}10^{-1}$ & $9.9{\times}10^{-2}$ & $6.4{\times}10^{-1}$ & $2.8{\times}10^{-1}$ & $2.7{\times}10^{-1}$ \\
64   & $3.0{\times}10^{-7}$ & $1.5{\times}10^{-4}$  & $1.8{\times}10^{-4}$ & $1.4{\times}10^{-2}$ & $1.5{\times}10^{-3}$ & $4.9{\times}10^{-2}$ & $9.3{\times}10^{-3}$ & $2.1{\times}10^{-1}$ \\
128  & $1.3{\times}10^{-12}$& $1.5{\times}10^{-11}$ & $5.3{\times}10^{-10}$& $3.9{\times}10^{-6}$ & $2.3{\times}10^{-7}$ & $6.5{\times}10^{-4}$ & $4.7{\times}10^{-5}$ & $2.4{\times}10^{-2}$ \\
256  & $1.7{\times}10^{-12}$& $3.6{\times}10^{-14}$ & $5.0{\times}10^{-13}$& $1.4{\times}10^{-12}$& $5.0{\times}10^{-13}$& $1.5{\times}10^{-8}$ & $1.5{\times}10^{-10}$& $6.5{\times}10^{-6}$ \\
512  & $1.9{\times}10^{-12}$& $2.5{\times}10^{-14}$ & $5.8{\times}10^{-13}$& $7.2{\times}10^{-13}$& $5.9{\times}10^{-13}$& $4.2{\times}10^{-11}$ & $4.9{\times}10^{-13}$& $5.9{\times}10^{-8}$ \\
\bottomrule
\end{tabular}%
}
\end{table}

Three findings emerge. First, erfc-trap converges to ${\sim}10^{-12}$
relative error by $Q = 128$ on uniform and adversarial stacks, and by
$Q = 256$ on the more demanding geometric and extreme distributions. At
the $Q = 256$ production default, all distributions achieve relative error
below $1.6 \times 10^{-10}$.

Second, tanh-sinh achieves a lower error floor on uniform distributions
(${\sim}10^{-14}$ at $Q = 256$) but degrades catastrophically at extreme
ratios, reaching only $6.5 \times 10^{-6}$ on the $10^9$ extreme case and
failing to improve further at $Q = 1024$. The double-exponential
substitution's strip of analyticity narrows as $S_{\max}/S_{\min}$ grows;
beyond some threshold, convergence fails entirely.

Third, this failure mode drove the choice of quadrature scheme. In the
typical case, tanh-sinh is the better method. But a general-purpose ICM
library must handle arbitrary stack distributions, including cases where
one player holds $10^9$ times as many chips as another. Gaussian
quadrature's uniform reliability - it keeps converging, predictably, all
the way to the $10^9$ bound - was judged more valuable than better
average-case precision.

\begin{figure}[t]
  \centering
  \includegraphics[width=0.7\textwidth]{../ICM/results/accuracy_convergence.png}
  \caption{Convergence of erfc-trap (Gauss-Legendre) quadrature toward
    machine precision as $Q$ increases, across the four stack
    distributions of Table~\ref{tab:accuracy}.}
  \label{fig:accuracy}
\end{figure}

\subsection{Calibration Constants}

Table~\ref{tab:constants} summarizes the constants that actually drive the
shipped dispatch and cost logic, all measured by direct microbenchmark or
offline profiling rather than fit by regression (Section~\ref{sec:calib-methodology}).
$\texttt{WRAP\_FMA\_NS}$ prices the tree engine's closed-form wrap-correction
cost; $\texttt{PAIRED\_CACHED\_CORR\_RATIO}$ and $\texttt{INDEP\_PAIR\_RATIO}$
scale measured FFT time to price the paired- and independent-correlate
phases; $L_2$ sizes the linear engine's checkpointing interval
(Section~\ref{sec:linear}); and the B-selection table size is the number
of empirically measured $(n,k,B)$ points behind
\texttt{select\_best\_B}$(n,k)$ (Section~\ref{sec:dispatch}). An earlier
roofline cost model built from streaming-bandwidth constants
($\tau_{\text{FMA}}$, $W_{L2}$, $W_{L3}$, $W_{\text{DRAM}}$) was replaced
by this direct-measurement approach and is no longer part of the shipped
library (Section~\ref{sec:calib-methodology}).

\begin{table}[t]
\centering
\caption{Platform calibration constants live in \texttt{libicm.a}'s dispatch path.}
\label{tab:constants}
\begin{tabular}{@{}lrrl@{}}
\toprule
\textbf{Constant} & \textbf{M3 Pro} & \textbf{Zen 4} & \textbf{Description} \\
\midrule
$\texttt{WRAP\_FMA\_NS}$ (ns) & 0.516 & 0.436 & Wrap-correction FMA cost \\
$\texttt{PAIRED\_CACHED\_CORR\_RATIO}$ & 1.560 & 1.547 & Paired cached correlate / full pipeline \\
$\texttt{INDEP\_PAIR\_RATIO}$ & 2.440 & 2.440 & \texttt{correlate\_fft\_pair} / full pipeline \\
$L_2$ (bytes) & 1\,MB & 1\,MB & Per-core L2, sizes checkpointing \\
B-selection table & 2{,}563 pts & 1{,}944 pts & Empirical $(n,k,B)$ lookup size \\
Typical $B$ & 32 & 32 & Most common empirical block size \\
FFT backend & FFTW+vDSP & AOCL (sole) & Library dispatch \\
\bottomrule
\end{tabular}
\end{table}

\subsection{Calibration Methodology: Direct Measurement vs.\ Indirect Regression}\label{sec:calib-methodology}

The empirical crossover table (Section~\ref{sec:dispatch}) replaced an
earlier roofline cost model that depended on hardware
constants: $\tau_{\text{FMA}}$, cache bandwidths, per-size FFT costs,
and several micro-architectural parameters; that are fit to timing data
by nonlinear least squares. Two of these constants proved structurally
unrecoverable from aggregate plan-timing data alone, leading to a
methodological shift from indirect regression to isolated direct
microbenchmarking. This section describes the failure mode, the fix, and
the precedent it follows.

\paragraph{Precedent: direct measurement over indirect inference.}
The principle of measuring a parameter directly rather than backing it out
of aggregate data has strong precedent in high-performance computing.
FFTW's \texttt{PATIENT}-mode planner measures the actual execution time of
candidate FFT plans directly rather than relying on an indirect cost
model~\citep{frigo2005}. The BeBOP/Sparsity system's register-blocking search empirically times
generated kernel variants per architecture to select the
fastest~\citep{demmel2004}. The
roofline model measures peak FLOP/s and memory bandwidth via dedicated
streaming microbenchmarks rather than inferring them from application
runs~\citep{williams2009}. In each case, the rationale is the same: when a
parameter's contribution to the total observed cost is small relative to
other, larger effects, an indirect regression lacks the signal to identify
it reliably.

\paragraph{Persistency of excitation.}
This failure mode has a formal name in control theory and system
identification: \emph{persistency of excitation}~\citep{ljung1999}. A
parameter is identifiable from training data only if the input signal
varies its effect enough to be distinguishable from noise and from the
effects of other parameters. When one parameter's contribution never
exceeds a few percent of any training example's total cost (as is the case
for wrap-correction cost in the tree engine; see below), the regression
objective is effectively flat with respect to that parameter: wildly
different values produce nearly identical aggregate predictions, and the
fitted value reflects optimizer drift rather than physical reality.

\paragraph{The wrap-correction identifiability failure.}
The tree engine's wrap correction (Algorithm~\ref{alg:tree}) corrects
aliasing wraparound when truncated polynomial multiplication is performed
via FFT with a transform size smaller than the full linear convolution
length. The cost of this correction, parameterized by
$\texttt{WRAP\_FMA\_NS}$ (nanoseconds per fused multiply-add in the
correction loop), is a small fraction of any sampled plan's total
predicted time. On AMD Zen~4, only 5 of 200 sampled plans had
wrap-correction cost exceeding 1\% of that plan's total predicted time
(maximum 1.5\%), making the constant effectively unidentifiable from
aggregate data.

The fix was to measure $\texttt{WRAP\_FMA\_NS}$ directly: an isolated
microbenchmark (\texttt{tools/bench\_wrap\_fma.c}) that executes only the
wrap-correction inner loop (a verbatim copy of the production code in
\texttt{src/icm.c}), sweeping the wrap range over $[4, 2048]$ across three
polynomial-size regimes (2K, 8K, 32K elements) so that the correction
dominates the measured time by construction ($R^2 = 0.9998$). The result,
displayed in Figure~\ref{fig:wrap-curve}, is not a single constant but a
genuine cache-hierarchy-transition curve: the marginal cost rises smoothly
from ${\sim}0.26$\,ns/FMA at small wrap ranges (L1/L2-resident,
near-FMA-throughput-bound) to ${\sim}0.53$\,ns/FMA at large wrap ranges
(memory-latency-bound as the working set exceeds cache capacity). At the
decision-relevant operating range ($\texttt{wrap\_m} \approx 64$--$384$,
matching tree levels~7--9 for typical $(n,k)$ pairs), the local slope is
${\sim}0.36$--$0.45$\,ns/FMA.

\begin{figure}[t]
  \centering
  \includegraphics[width=0.65\textwidth]{../ICM/results/wrap_fma_cost_curve.png}
  \caption{Measured marginal cost of the wrap-correction FMA loop as a
    function of wrap range, for the SMALL\_2048 polynomial-size regime
    (Zen~4). The curvature reflects a real cache-hierarchy transition:
    L1/L2-resident at small $\texttt{wrap\_m}$ (${\sim}0.26$\,ns/FMA),
    transitioning to memory-latency-bound at large $\texttt{wrap\_m}$
    (${\sim}0.53$\,ns/FMA). The production value
    $\texttt{WRAP\_FMA\_NS} = 0.436$ (dashed line) is a representative
    point estimate from the least-squares slope over the
    $\texttt{wrap\_m} \in [64, 384]$ operating window.}
  \label{fig:wrap-curve}
\end{figure}

\paragraph{Disclosure: point-estimate approximation.}
The shipped $\texttt{WRAP\_FMA\_NS}$ values, $0.436$ on Zen~4 and
$0.516$ on Apple M3~Pro, are representative point estimates from a
hand-chosen operating-range window ($\texttt{wrap\_m} \in [64, 384]$),
not a fully piecewise or continuous model of the cache-hierarchy
transition visible in Figure~\ref{fig:wrap-curve}. The dispatch model
treats wrap correction as having a single constant marginal cost; a more
accurate model would account for the dependence on working-set size, but
the single-constant approximation has proved sufficient in practice
(dispatch accuracy 42/44 on Zen~4, Section~\ref{sec:dispatch}).

\paragraph{M3 Pro: a second identifiability failure.}
The same methodology was applied to $\texttt{FP64\_DIV\_NS}$ on the Apple
M3~Pro after the indirect fit converged to a physically implausible value
(0.5\,ns, at its imposed lower bound, well below the latency a
dependency-chained FP64 division is expected to take on this
architecture). A second dedicated microbenchmark
(\texttt{tools/bench\_div\_chain.c}) measured the dependency-chained
division cost directly (matching the actual usage in the hybrid engine's
leaf-extraction synthetic-division recurrence, Section~\ref{sec:hybrid}),
yielding $3.4890$\,ns.

\paragraph{Collinearity caveat (resolved).}
An earlier revision of this calibration pinned only
$\texttt{WRAP\_FMA\_NS}$ and $\texttt{FP64\_DIV\_NS}$ directly, leaving the
remaining cost-model constants to a nonlinear-least-squares fit. That fit
raised the aggregate RMS log-relative error from 6.57\% to 10.2\% on
M3~Pro and pushed $\texttt{FFT\_OVERHEAD\_NS}$ to a physically odd
631\,ns to compensate, a collinearity artifact of the
\texttt{sample\_plans} training data, not a correctness issue.
\textbf{This limitation no longer applies.} All six scalar constants are
now pinned directly from microbenchmarks; \texttt{tools/fit\_cost\_model.py}
skips the regression entirely once every scalar is supplied, and
$\texttt{FFT\_OVERHEAD\_NS}$ is $0.0$ on every device. The identifiability
failure recorded here is what motivated that migration, not a standing
limitation of anything shipped. \texttt{bench\_grid verify} passes all
tests and \texttt{bench\_grid crossover} shows a clean, monotonic
linear-to-hybrid transition at $k \approx 122$--$124$ on M3~Pro and
$k \approx 249$--$281$ on Zen~4.

\subsection{Subset Equity Pruning}\label{sec:subset}

The hybrid engine supports a subset-query mode that computes equities for
only a specified subset of target players. It exploits
\emph{target-locality pruning}: a per-block hot/cold bitmask is propagated
bottom-up through the tree, and cold branches are skipped during the
top-down propagation phase. This optimization is correct (verified
bit-exact against the full-equity engine) but has a fixed upfront cost for
mask construction and per-node branch checks.

Table~\ref{tab:subset} reports speedups measured on M3~Pro. The pruning
delivers a net win across the tested range of $n$ at low target fractions,
topping out at ${\sim}1.5\times$ speedup at $n = 1024$ with 1\% targets, and
fading toward breakeven as the target fraction approaches 100\%.

\begin{table}[t]
\centering
\caption{Speedup of subset equity vs.\ full-equity computation, M3~Pro,
  single-threaded, $Q = 256$. Ratios ${<}\,1$ indicate net slowdown.}
\label{tab:subset}
\begin{tabular}{@{}l c c c@{}}
\toprule
$n$ & Speedup @1\% targets & Speedup @5--10\% targets & Speedup @$\ge$25\% targets \\
\midrule
1024  & 1.49 & 1.17--1.21 & 1.00--1.15 \\
4096  & 1.24 & 1.11--1.13 & 0.99--1.09 \\
16384 & 1.16 & 1.10 & 1.00--1.08 \\
65536 & 1.09 & 1.05 & 1.00--1.03 \\
\bottomrule
\end{tabular}
\end{table}

% ======================================================================
\section{GPU Implementation: NVIDIA B200}\label{sec:gpu}
% ======================================================================

The GPU implementation targets the NVIDIA B200 and adapts the hybrid
block-tree algorithm to the GPU execution model. The linear engine is
retained in the dispatch code only as a trivial-case fallback ($n < 16$ or
$k_{\text{pad}} < 4$); its sequential player-by-player structure cannot
saturate GPU parallelism at any size that matters in practice, and it is
never selected across the full calibrated heatmap (210/210 cells choose
hybrid). Only the tree-based engines, where all nodes at each level are
independent, map naturally to GPU execution at scale.

\subsection{Architecture}

The tree's level-by-level structure becomes a sequence of batched kernel
launches. At the bottom levels of a tree with $n = 1{,}572{,}864$ leaves
and $Q = 256$, this yields millions of independent operations - far more
than needed to saturate the GPU's streaming multiprocessors.

\subsection{Three-Tier Kernel Dispatch}

The planner assigns each subproduct-tree level to one of three kernel
tiers based on polynomial degree:
\begin{itemize}
  \item \textbf{Schoolbook} ($d < 64$): shared-memory kernels with
    excellent coalescing and zero launch overhead.
  \item \textbf{cuFFTDx fused} ($64 \le d \le 4096$): device-side FFT
    kernels that fuse forward R2C, pointwise multiply, and inverse C2R
    into a single on-device launch, eliminating intermediate global-memory
    round-trips.
  \item \textbf{Batched cuFFT} ($d > 4096$): globally-optimized plans
    for the highest-throughput regime.
\end{itemize}

The entire multi-level execution sequence, spanning dozens of kernel launches
across all three tiers - is captured into a single CUDA Graph and replayed
with negligible launch overhead.

\subsection{Key Optimizations}

\paragraph{FFT-stride layout.}
The single largest optimization: polynomials are stored at their
FFT-padded stride rather than at their logical degree. This eliminates
the gather/scatter passes that cuFFT's batched interface would otherwise
require, cutting two full memory passes per FFT call.

\paragraph{Arena allocator.}
The baseline implementation made 35 separate \texttt{cudaMalloc} calls.
Replacing these with a single arena allocation eliminates serialization
at the CUDA driver layer.

\paragraph{Q-batch pipelining.}
Rather than processing all $Q = 256$ quadrature points in one monolithic
batch, I process $Q_{\text{batch}} = 4$--$8$ points at a time, overlapping
the block-build phase of one batch with tree propagation of the previous
batch using multi-stream execution.

\paragraph{Rejected optimizations.}
Two approaches were tested and rejected after benchmarking:
\begin{itemize}
  \item \textbf{cuFFT LTO callbacks:} Fusing the pointwise multiply into
    cuFFT's load/store callbacks was rejected on accuracy grounds, not
    speed: per-element callback rounding compounds across tree levels to
    ${\sim}2\%$ error, despite being within cuFFT's own accuracy
    specification -- unacceptable given this algorithm's FP64 requirement
    (Section~\ref{sec:experiments}).
  \item \textbf{$B = 1$ balanced tree:} Using single-player leaves (no
    block structure) was substantially slower because the tree performs
    more FMA operations and adds kernel launch barriers.
\end{itemize}

\subsection{GPU Calibration Methodology}\label{sec:gpu-calib-methodology}

Block-size selection on GPU follows the same empirical-lookup principle as
the CPU (Section~\ref{sec:dispatch}): \texttt{gpu\_empirical\_best\_B}$(n,k)$
performs a joint $(n,k)$ nearest-neighbor lookup over a calibrated
\texttt{gbselect\_n[]}/\texttt{gbselect\_k[]}/\texttt{gbselect\_B[]} table, in
log space, exactly mirroring the CPU's \texttt{select\_best\_B}. An
analytical fallback (summed FMA/FFT cost estimation) remains in
\texttt{gpu\_select\_best\_B\_est} for the case where no calibrated candidate
satisfies the $(n,k)$ bounds at all; in normal operation this path is not
exercised. Engine dispatch (linear vs.\ hybrid) is the one GPU decision that
still compares candidates analytically rather than from a measured table, but
it never selects linear across the full 210-cell calibrated heatmap --
consistent with Section~\ref{sec:gpu}'s observation that the linear engine's
sequential structure cannot compete with the tree engines' parallelism at any
GPU-relevant size.

\paragraph{Deliberately targeted scope, not a full sweep.}
Unlike the CPU tables, the GPU B-selection table above the $k = n$ frontier
($n > 1{,}572{,}864$) is not built from a dense adaptive sweep. A full
48-candidate sweep to $n = 33{,}554{,}432$ was measured, from the
calibration heatmap itself, to cost roughly 4.5 hours and \$28 of rented
B200 time, with $n = 33{,}554{,}432$ alone accounting for 54\% of total
heatmap compute. Instead, 13 targeted anchor points
($n \in \{2\text{M}, 4\text{M}, 8.4\text{M}\}$ at $k = n/8, n/4, n/2, n$, plus
$n = 16.7\text{M}$ at $k = n$ only) were measured directly with
\texttt{calibrate\_gpu\_best\_b.cu --narrow-around}; all 13 converged
uniformly on $B = 128$. The full adaptive calibration tooling
(\texttt{calibrate\_adaptive.py}) works and is available for third parties
porting to new hardware; only this project's own above-frontier B200 data is
a deliberate targeted subset rather than the CPU's full treatment.

\paragraph{Bug 1: single-repetition timing noise.}
An attempt to extend the GPU calibration into the previously under-sampled
$n = 1{,}024$--$524{,}288$ region, using the same adaptive priority-queue
refinement already proven on CPU (Section~\ref{sec:dispatch}), initially
produced a table that was measurably worse than the one it replaced: a diff
against the pre-run baseline over the standard 84-cell benchmark grid showed
36 regressions, several severe (e.g.\ $n = 2{,}048, k = 512$: dispatched
$B = 1{,}280$ against a true $B = 64$, a 292.75\% gap; $n = 16{,}384, k = 64$:
dispatched $B = 1{,}024$ against a true $B = 48$, a 265.32\% gap), and total
grid time \emph{increased} 4.57\% relative to the unmodified baseline. Both
worst cells finished in under $\sim$2\,ms total -- GPU kernel-launch-overhead
territory. The root cause: the measurement oracles
(\texttt{validate\_planner\_gpu.cu}, \texttt{calibrate\_gpu\_best\_b.cu}, and
their CPU counterparts \texttt{validate\_best\_b.c}/
\texttt{calibrate\_best\_b.c}) timed each candidate $B$ with a single
repetition and no noise check, taking the fastest single sample as the
winner. The GPU's candidate range spans $B = 16$ to $B = 1{,}536$
(96$\times$) versus the CPU's $\{8,16,24,32,48,64\}$ (8$\times$), so a
noise-driven misranking at small, sub-few-millisecond problem sizes could
select a structurally different, far worse block size once dispatched for
real -- a failure mode the CPU's narrower candidate range and larger
per-measurement amortization (256 internal queries per probe) largely
avoids. The fix ported the adaptive median/coefficient-of-variation
convergence loop already proven in \texttt{heatmap\_gpu.cu} (up to 10
repetitions, extended to 15 if $\mathrm{cv} > 3\%$) into all four
measurement tools, GPU and CPU. A direct re-probe of the two worst cells
after the fix confirmed the diagnosis: the $n = 2{,}048, k = 512$ gap
collapsed from 292.75\% to $-0.10\%$, and $n = 16{,}384, k = 64$ from
265.32\% to 0.34\%, while a third, independently real regression at
$n = 65{,}536, k = 2{,}048$ stayed at 17.76\% (previously documented at
16.5\%), confirming the fix removes fake noise-driven signal without
masking genuine ones. The same night, a real, full-domain M3~Pro CPU
recalibration through the fixed tool (92 probes, 2{,}466 $\to$ 2{,}558
points) showed only modest gaps throughout (0--8\%, no wild swings),
consistent with the CPU's narrower candidate range being far less exposed to
this mechanism.

\paragraph{Bug 2: sparse nearest-neighbor coverage.}
With the measurement noise fixed and independently confirmed, a second,
full attempt to populate the same $n = 1{,}024$--$524{,}288$ region (252
probes, 93 $\to$ 345 points) surfaced a second, unrelated failure: a full
diff against the original baseline showed 40 regressions out of the same
84-cell grid, several catastrophic ($n = 1{,}024, k = 128$: $B = 64 \to
1{,}024$, $+402.7\%$; $n = 65{,}536, k = 128$: $B = 64 \to 1{,}536$,
$+329.3\%$; $n = 32{,}768, k = 256$: $B = 64 \to 1{,}536$, $+275.1\%$), and
total grid time 5.23\% worse than the original baseline -- worse than
Bug~1's own regression despite every individual measurement now being
independently verified correct. This is a distinct mechanism from Bug~1:
correct individual measurements do not by themselves make a sparse
nearest-neighbor table safe. $B^*(n,k)$ has real local structure in this
region, unlike the large-$n$ region the targeted anchors above already
cover; populating a previously empty region from scratch via priority-queue
sampling makes each newly injected point the nearest neighbor for an entire
neighborhood of queries that were never themselves directly probed, with
nothing to check whether that neighborhood actually agrees with the one
point that was measured. This table was never shipped: it existed only on
the rented instance used to produce it and was discarded before the
instance was destroyed, preserved separately as evidence
(\texttt{results/gpu\_fft\_config\_20260731\_run2\_BROKEN\_evidence.h}), not
as usable calibration data.

\paragraph{Canary safety net.}
In direct response to both failures, \texttt{calibrate\_adaptive.py} now
includes an automatic safety net rather than relying on a human diff after
the fact. Before any calibration run, it snapshots the target config header
byte-for-byte and measures a small, fixed, held-out canary set (three
$k$-fractions per landmark $n$-anchor, deliberately disjoint from the
landmark and adaptive-refinement points) using the same converged
measurement oracle the rest of the pipeline trusts. After the run, it
rebuilds the validation binary against the caller-supplied
\texttt{--rebuild-cmd} (required; no default is assumed safe, since the GPU
build needs an environment-specific include path and the CPU validator has
no default binary at all) and re-measures the identical canaries,
auto-reverting to the exact pre-run byte-for-byte snapshot if the aggregate
canary time regresses more than 2\% or any single canary regresses more
than 5\%. The mechanism was validated three ways before being trusted: a
live pass-path run confirmed byte-identical (via \texttt{md5sum}) when
canaries agree; a live run with a deliberately broken
\texttt{--rebuild-cmd} confirmed the file is restored byte-for-byte; and
unit tests replaying this session's own worst regression
($n = 1{,}024, k = 128$, $+402.7\%$) confirm the comparison logic fails in
exactly the way it should.

\paragraph{Disclosure: an open limitation, not retroactively closed.}
The canary net makes the \emph{next} calibration attempt over the
$n = 1{,}024$--$524{,}288$ region safe to run unattended -- it does not
itself produce the denser sampling that region would need for the
nearest-neighbor lookup to be locally valid there. As of this writing, that
region of the GPU B-selection table is exactly as calibrated as it was
before this investigation: neither worse (both broken attempts were
discarded before shipping) nor better. This is disclosed here rather than
silently left implicit, consistent with this paper's broader position that
the GPU calibration surface received a narrower, more targeted treatment
than the CPU's, not the same full adaptive sweep.

% ======================================================================
\section{GPU Experimental Results}\label{sec:gpu-results}
% ======================================================================

\textbf{Hardware.} NVIDIA B200, sm\_100, ${\sim}$180\,GB HBM3e. All
experiments use $Q = 256$ quadrature points and cuFFTDx fused device-side
kernels with CUDA graph capture (Section~\ref{sec:gpu}).

Table~\ref{tab:gpu} reports a systematic $(n,k)$ grid sampled directly from
the same 211-point calibration heatmap that generates
Figure~\ref{fig:gpu-contour}, in the same row/column layout as
Tables~\ref{tab:serial}--\ref{tab:parallel} so the CPU and GPU results are
directly comparable. Table~\ref{tab:gpu-frontier} separately reports the
one-second threshold located by a dedicated binary search on measured time
(\texttt{tools/threshold\_search\_gpu.cu}, median of 5 repetitions per
candidate $n$) rather than by interpolating the systematic grid - these are
the specific $n$ values that pin down the 1-second boundaries quoted below,
and are not a subsample of Table~\ref{tab:gpu}.

\begin{table}[t]
\centering
\caption{NVIDIA B200 GPU execution time (ms), $Q = 256$, FP64. Systematic
  $(n,k)$ grid sampled from the calibration heatmap.}
\label{tab:gpu}
\begin{tabular}{@{}l rrrr@{}}
\toprule
$n$ & $k{=}64$ & $k{=}1024$ & $k{=}n/2$ & $k{=}n$ \\
\midrule
4{,}096        & 0.37     & 0.76     & 0.87      & 0.89 \\
16{,}384       & 1.18     & 2.81     & 3.93      & 4.17 \\
65{,}536       & 4.29     & 9.46     & 19.43     & 20.26 \\
262{,}144      & 16.58    & 36.38    & 95.95     & 100.09 \\
1{,}048{,}576  & 65.68    & 178.26   & 504.29    & 509.51 \\
4{,}194{,}304  & 272.47   & 671.06   & 2{,}352.43 & 2{,}320.45 \\
33{,}554{,}432 & 2{,}506.71 & 5{,}059.26 & 22{,}321.49 & 22{,}865.76 \\
\bottomrule
\end{tabular}
\end{table}

\begin{table}[t]
\centering
\caption{NVIDIA B200 one-second threshold, located by binary search on
  measured time rather than sampled from the grid above. Each pair brackets
  the 1000\,ms crossing; the search terminated at a bracket width of 15{,}360
  on both curves.}
\label{tab:gpu-frontier}
\begin{tabular}{@{}l rr@{}}
\toprule
Curve & $n$ & Time (ms) \\
\midrule
$k = n$   & 1{,}490{,}944 & 918.8 \\
          & 1{,}506{,}304 & 1{,}124.2 \\
\addlinespace
$k = 100$ & 7{,}975{,}936 & 998.9 \\
          & 7{,}991{,}296 & 1{,}001.1 \\
\bottomrule
\end{tabular}
\end{table}

\begin{figure}[t]
  \centering
  \includegraphics[width=0.6\textwidth]{../ICM/results/gpu_contour.png}\\[1em]
  \includegraphics[width=0.75\textwidth]{../ICM/results/runtime_vs_n_gpu.png}
  \caption{NVIDIA B200. Top: one-second contour (maximum $n$ achievable
    within 1\,s as a function of $k$), the GPU analogue of
    Figure~\ref{fig:contour}. Bottom: runtime vs.\ $n$ at fixed $k$,
    showing the tree-level cliffs referenced below.}
  \label{fig:gpu-contour}
\end{figure}

At $k = n$, the GPU computes all 262{,}144 player equities in 100.1\,ms.
The under-one-second frontier for $k = n$ is $n = 1{,}490{,}944$
(918.8\,ms); the next candidate the search probed above it,
$n = 1{,}506{,}304$, crosses to 1{,}124.2\,ms. At $k = 100$ the frontier
reaches $n = 7{,}975{,}936$ players (998.9\,ms). The frontier is limited by
VRAM at large $k$ and by tree-level cliffs at large $n$.

% ======================================================================
\section{Conclusion}\label{sec:conclusion}
% ======================================================================

I have presented a high-performance algorithm for computing ICM tournament
equities using generating-function quadrature and FFT-accelerated
subproduct trees. The algorithm achieves $\bigO(Qn\log^2 k)$ time, double-precision accuracy via $Q = 256$ quadrature points (relative error
${<}\,5 \times 10^{-12}$ on typical inputs), and automatically dispatches
between a linear engine (optimal for small $k$) and a hybrid block-tree
engine (optimal for large $k$) using an empirically calibrated crossover
table (Section~\ref{sec:dispatch}).

The propagation phase is precisely the reverse-mode derivative of the
build phase, and recognizing this gives both a clean correctness proof
(Theorem~\ref{thm:tree-correct}) and complexity parity between build and
propagation for free. The dominant cost throughout is data movement, not
arithmetic: the empirical dispatch table, L2-aware checkpointing, and the
GPU's FFT-stride layout all exist because moving data between cache levels
is the binding constraint, not the FMA count. And the cost model and FFT
dispatch decisions throughout - per-size FFT timing, per-level cache
bandwidths, dual-library selection - are driven by direct calibration
against measured hardware behavior rather than predicted from first-principles
constants, because measuring turned out to be cheaper and more reliable
than modeling it a priori.

On the CPU, the batched linear and hybrid engines compute exact equities
for up to 17,984 players in one second, single-threaded, on AMD Zen~4. A
CUDA implementation for NVIDIA B200 extends this to roughly 1.5 million
players in about a second, enabling ICM analysis at field sizes that were
previously intractable.

Future work includes exploring real-time tournament decision support where
subset equity queries can exploit the active-player masking already present
in the implementation. A more speculative direction: since the propagation
phase is already reverse-mode automatic differentiation of the build phase
(Section~\ref{sec:tree}), and $a_j(v) = v^{S_j}$ is itself differentiable
in $S_j$, extending this AD lens one step further - differentiating through
the quadrature integral itself - would give exact gradients
$\partial E_i / \partial S_j$ at the same $\bigO(Qn\log^2 k)$ cost as the
equities themselves, potentially useful as a differentiable ICM node in
gradient-based tournament-strategy training objectives.

% ======================================================================
% APPENDIX (optional, but useful for the mathematically inclined reader)
% ======================================================================

\appendix
\section{Closed-Form Reference Equities}\label{sec:appendix}

The library is validated against exact closed-form reference values for two
special payout structures. Both follow from a single combinatorial identity
applied at different orders, combined with linearity of expectation - not
from enumerating elimination orderings - and are exact for any $n$.

Write player $i$'s finishing position as $m = M + 1$, where $M$ is the
(random) number of other players who finish ahead of $i$ ($M = 0$ is first
place). For any $t \ge 0$, the number of ways to choose $t$ of the
$n - 1 - M$ players who finish \emph{behind} $i$ satisfies, on every
realized outcome,
\begin{equation}\label{eq:combinatorial-identity}
  \binom{n-1-M}{t} = \sum_{\substack{T \subseteq [n]\setminus\{i\} \\ |T| = t}}
  \mathbf{1}[i \text{ finishes better than every player in } T],
\end{equation}
since this is exactly choosing $t$ of the $n-1-M$ players ranked below $i$ -
no probability is involved yet, it is a counting identity on one outcome.
Taking expectations and applying linearity term-by-term to the sum of
indicators on the right (valid regardless of how the indicator events are
correlated with each other):
\begin{equation}\label{eq:combinatorial-expectation}
  E\!\left[\binom{n-1-M}{t}\right] = \sum_{\substack{T \subseteq [n]\setminus\{i\} \\ |T|=t}}
  \Pr[i \text{ beats every player in } T].
\end{equation}
``$i$ beats every player in $T$'' means $i$'s exponential clock is the
smallest among $\{i\} \cup T$; by the same competing-exponentials argument
as Section~\ref{sec:derivation} (the minimum of independent
$\mathrm{Exp}(S_j)$ variables is itself exponential with rate equal to the
sum of the rates), $\Pr[i \text{ beats every player in } T] = S_i /
(S_i + \sum_{j \in T} S_j)$.

\paragraph{V1 (linear payout, $\pi_m = n - m + 1$).} Since $n - M =
\binom{n-1-M}{0} + \binom{n-1-M}{1}$, apply Equation~\ref{eq:combinatorial-expectation}
at $t=0$ (trivially $1$, the empty subset $T = \emptyset$) and $t=1$ (one
term per opponent $j$, $T = \{j\}$):
\begin{align*}
  E_i = E[n - M]
  &= E\!\left[\binom{n-1-M}{0}\right] + E\!\left[\binom{n-1-M}{1}\right] \\
  &= 1 + \sum_{j \neq i} \Pr[i \text{ beats every player in } \{j\}] \\
  &= 1 + \sum_{j \neq i} \frac{S_i}{S_i + S_j}.
\end{align*}

\paragraph{V2 (quadratic payout, $\pi_m = \binom{n-m}{2}$).} Since
$\pi_m = \binom{n-m}{2} = \binom{n-1-M}{2}$ directly, apply
Equation~\ref{eq:combinatorial-expectation} at $t=2$ (one term per opponent
pair $T = \{j,k\}$):
\begin{align*}
  E_i = E\!\left[\binom{n-1-M}{2}\right]
  &= \sum_{j < k,\, j,k \neq i} \Pr[i \text{ beats every player in } \{j,k\}] \\
  &= \sum_{j < k,\, j,k \neq i} \frac{S_i}{S_i + S_j + S_k}.
\end{align*}

Higher payout schedules follow the same pattern for larger $t$; V1 and V2
are the $t \le 2$ cases used here as exact, closed-form, arbitrary-$n$
ground truth. Both are implemented as \texttt{icm\_v1\_exact()} and
\texttt{icm\_v2\_exact()} in the reference library and used to validate the
quadrature-based implementation across all stack distributions.

\bibliographystyle{plainnat}
\bibliography{references}

\end{document}
